\documentclass[reqno]{amsart}

% Packages
\usepackage{import}
\usepackage{xcolor}
\usepackage{cancel}
\usepackage{mathtools}
\usepackage{hyperref}
\usepackage{setspace}
\usepackage{float}
\usepackage[shortlabels]{enumitem}
\usepackage[margin=1in]{geometry}
\setlength{\parindent}{0pt}
\usepackage{titlesec}
\titlespacing*{\section}{0pt}{0pt}{25pt}

\usepackage{scalerel}

\renewcommand{\vec}{\mathbf}
\DeclareMathOperator{\vectorprojection}{proj}
\NewDocumentCommand{\proj}{s e{_} m}{%
    \IfBooleanTF{#1}{#3_{\stretchrel*{\parallel}{\perp}#2}}{\vectorprojection_{#2}#3}
}
\renewcommand{\Re}{\operatorname{Re}}
\renewcommand{\Im}{\operatorname{Im}}
\renewcommand{\Arg}{\operatorname{Arg}}

% Custom perms and combs commands
\newcommand\Myperm[2][^n]{\prescript{#1\mkern-2.5mu}{}P_{#2}}
\newcommand\Mycomb[2]{\prescript{#1\mkern-0.5mu}{}C_{#2}}
% Custom column vector command
\newcommand{\cvec}[1]{ \begin{pmatrix} #1 \end{pmatrix} }
% Custom determinant command
\newcommand{\det}[1]{ \begin{vmatrix} #1 \end{vmatrix} }

% Colors
% \usepackage[T1]{fontenc}			% Font package
\usepackage{fouriernc}

\usepackage{sectsty}
\usepackage{graphicx}
\usepackage{amsmath}
\usepackage{amsfonts}
\usepackage{amssymb}

\usepackage{tikz}
\usepackage{eso-pic}
\usetikzlibrary{calc, shadows.blur}
\usetikzlibrary{angles, quotes}
\usetikzlibrary{3d}

\setlength{\jot}{1.5em}
\setlength{\thickmuskip}{10}
\setlength{\medmuskip}{7}
\setlength{\thinmuskip}{5}
\setstretch{1.35}

\title{%
	MATH1241: Higher Mathematics 1B \\
	Assignment (Applied Mathematics Flavour)}
\author{L. Cheung (z5763342)}
\date{July 14, 2026}

\begin{document}

\maketitle

\section*{\LARGE{Question 1}}

	The function $T: \mathbb{R}^3 \to \mathbb{R}^2$ is defined by
	\begin{align*}
		T\cvec{x_1\\x_2\\x_3} = \cvec{3x_2 - 2x_3\\-2x_1 - 4x_3} \quad \text{for all } \cvec{x_1\\x_2\\x_3} \in \mathbb{R}^3.
	\end{align*}
	Show that $T$ is linear.
	\\ \\
	The function $T$ is a linear map if it preserves both vector addition and scalar multiplication. Let $\vec{u}$ and $\vec{v}$ be vectors in $\mathbb{R}^3$,
	\begin{align*}
		T(\vec{u} + \vec{v}) &= T\cvec{u_1 + v_1\\u_2 + v_2\\u_3 + v_3}	\\
							 &= \cvec{3(u_2 + v_2) - 2(u_3 + v_3) \\ -2(u_1 + v_1) - 4(u_3 + v_3)}	\\
							 &= \cvec{3u_2 + 3v_2 - 2u_3 - 2v_3 \\ -2u_1 -2v_1 - 4u_3 - 4v_3}	\\
							 &= \cvec{3u_2 - 2u_3 \\ -2u_1 - 4u_3} + \cvec{3v_2 - 2v_3 \\ -2v_1 - 4v_3}	\\
							 &= T(\vec{u}) + T(\vec{v}).
	\end{align*}
	Since $T(\vec{u} + \vec{v}) = T(\vec{u}) + T(\vec{v})$, the map $T$ preserves vector addition. Furthermore, let $\lambda \in \mathbb{F}$,
	\begin{align*}
		T(\lambda \vec{u}) &= T\cvec{\lambda u_1\\\lambda u_2\\\lambda u_3}	\\
						   &= \cvec{3(\lambda u_2) - 2(\lambda u_3) \\ -2(\lambda u_1) - 4(\lambda u_3)}	\\
						   &= \cvec{3\lambda u_2 - 2\lambda u_3 \\ -2\lambda u_1 - 4\lambda u_3}	\\
						   &= \cvec{\lambda (3u_2 - 2u_3) \\ \lambda (-2u_1 - 4u_3)}	\\
						   &= \lambda \cvec{3u_2 - 2u_3 \\-2u_1 - 4u_3}	\\
						   &= \lambda T(\vec{u}).
	\end{align*}
	Since $T(\lambda \vec{u}) = \lambda T(\vec{u})$, the function $T$ preserves scalar multiplication. The mapping $T$ passes both conditions of a linear transformation, hence $T$ is linear.

\newpage
\section*{\LARGE{Question 2}}

	Show that
	\begin{align*}
		S = \left\{ x \in \mathbb{R}^4 : 4x_1 - 6x_2 + 7x_3 - 9x_4 = 0 \right\}
	\end{align*}
	is a subspace of $\mathbb{R}^4$.
	\\ \\
	A subset $S$ of a vector space $V$ over $\mathbb{F}$ is a subspace if and only if $S$ contains the zero vector of the vector space $V$, $S$ is closed under vector addition, and $S$ is closed under scalar multiplication.

	The given subset $S$ contains the zero vector of $\mathbb{R}^4$, $\vec{x} = \cvec{0\\0\\0\\0}$. Consider two vectors in the subset $S$, $\vec{u}$ and $\vec{v}$. By definition, each vector satisfies
	\begin{align}
		\tag{i}
		4u_1 - 6u_2 + 7u_3 - 9u_4 = 0,		\\
		\tag{ii}
		4v_1 - 6v_2 + 7v_3 - 9v_4 = 0,
	\end{align}
	respectively. Substituting the sum of the two vectors, $\vec{u} + \vec{v}$, into the defining expression yields
	\begin{align*}
		4(u_1 + v_1) - 6(u_2 + v_2) + 7(u_3 + v_3) - 9(u_4 + v_4) &= (4u_1 - 6u_2 + 7u_3 - 9u_4) + (4v_1 - 6v_2 + 7v_3 - 9v_4)	\\
																  &= 0 + 0 \quad \text{by (i) and (ii)}	\\
																  &= 0.
	\end{align*}
	Since $\vec{u} + \vec{v}$ satisfies the defining equation of $S$, the subset is closed under vector addition. Furthermore, let $\vec{u}$ be an arbitrary vector in $S$, and $c$ be an arbitrary scalar in $\mathbb{R}$. Consider the scalar multiple of $\vec{u}$ by c,
	\begin{align*}
		c\vec{u} &= \cvec{c u_1 \\ c u_2 \\ c u_3 \\ c u_4}.
	\end{align*}
	This resulting vector must be contained within $S$, that is
	\begin{align*}
		4cu_1 - 6cu_2 + 7cu_3 - 9cu_4 &= c \cdot (4u_1 - 6u_2 + 7u_3 - 9u_4)	\\
				  					  &= c \cdot 0 \quad \text{from (i)}	\\
				  					  &= 0.
	\end{align*}
	Since $S$ contains the zero vector, is closed under vector addition, and is closed under scalar multiplication, the Subspace Theorem states that it is a subspace of $\mathbb{R}^4$.

\newpage
\section*{\LARGE{Question 3}}

	Let $\mathbb{P}_n$ be the set of real polynomials of degree at most $n$. Show that
	\begin{align*}
		S = \left\{P \in \mathbb{P}_7: p(-6) = p(7) \right\}
	\end{align*}
	is a subspace of $\mathbb{P}_7$.
	\\ \\
	By the Subspace Theorem, the set $S$ is a subspace of $V$ if and only if it contains the zero vector of $V$, $S$ is closed under vector addition, and closed under scalar multiplication.
	\\ \\	
	Let $\vec{0}$ denote the zero polynomial $\vec{0}(x) = 0 \; \forall x \in \mathbb{R}$. Therefore we have
	\begin{align*}
		\vec{0}(-6) = 0 \quad \text{and} \quad \vec{0}(7) = 0.
	\end{align*}
	Since $\vec{0}(-6) = \vec{0}(7)$, the zero polynomial satisfies the defining condition of $S$, so $\vec{0} \in S$.
	\\ \\
	Let $p,q \in S$ be arbitrary polynomials. By definition,
	\begin{align*}
		\tag{i}
		p(-6) = p(7),	\\
		\tag{ii}
		q(-6) = q(7).
	\end{align*}
	The sum $p + q$ is contained in the vector space $\mathbb{P}_7$ since the sum of two polynomials of degree at most 7 results in a polynomial of degree at most 7. Evaluating $p + q$ at $-6$ and $7$ gives
	\begin{align*}
		(p + q)(-6) &= p(-6) + q(-6)		\\
					&= p(7) + q(7) \quad \text{by (i) and (ii)}			\\
					&= (p + q)(7).
	\end{align*}
	Since $(p + q)(-6) = (p + q)(7)$, the set $S$ is closed under addition.
	\\ \\
	Finally, let $p \in S$ be arbitrary such that
	\begin{align*}
		\tag{iii}
		p(-6) = p(7),
	\end{align*}
	and let $c$ be an arbitrary scalar. The scalar multiple $cp \in \mathbb{P}_7$ since a scalar multiple of a polynomial of degree at most 7 is a polynomial of degree at most 7. Evaluating $cp$ at $-6$ and $7$ gives
	\begin{align*}
		(cp)(-6) &= c \cdot p(-6)	\\
				 &= c \cdot p(7) \quad \text{by (iii)}	\\
				 &= (cp)(7).
	\end{align*}
	Since $(cp)(-6) = (cp)(7)$, $cp \in S$ and the subset $S$ is closed under scalar multiplication. Since the subset of $\mathbb{P}_7$, $S$, contains the zero vector, is closed under vector addition, and closed under scalar multiplication, it is a subspace of $\mathbb{P}_7$.

\newpage
\section*{\LARGE{Question 4}}

	The air in a 84 cubic metre kitchen is initially clean, but when Ian burns his toast while making breakfast, smoke is mixed with the room's air at a rate of 0.05 mg per second. An air conditioning system exchanges the mixture of air and smoke with clean air at a rate of 9 cubic metres per minute. Assume that the pollutant is mixed uniformly throughout the room and that burnt toast is taken outside after 72 seconds. Let $S(t)$ be the amount of smoke in mg in the room at time $t$ (in seconds) after the toast first began to burn.

	\begin{enumerate}[(a)]
		\item Find a differential equation obeyed by $S(t)$.
			\\ \\
			Smoke enters the room at a constant rate of 0.05 mg and leaves at a rate proportional to the current concentration, giving the equation
			\begin{align*}
				\frac{dS}{dt} = 0.05 - \frac{0.15}{84} \cdot S = 0.05 - \frac{S}{560} \quad \text{for } 0 \leq t \leq 72.
			\end{align*}
			\\
		\item Find $S(t)$ for $0 \leq t \leq 72$ by solving the differential equation in (a) with an appropriate initial condition.
			\\ \\
			The differential equation is a first-order separable ODE and can be solved using an integrating factor. Rearranging the equation gives
			\begin{align*}
				\frac{dS}{dt} + \frac{1}{560} \cdot S = 0.05.
			\end{align*}
			The integrating factor can be calculated as $e^{\int \frac{1}{560} dt} = e^{\frac{t}{560}}$. Multiplying both sides of the differential equation by the integrating factor gives
			\begin{align*}
				e^{\frac{t}{560}} \left[\frac{dS}{dt} + \frac{1}{560} \cdot S\right] &= e^{\frac{t}{560}} \cdot 0.05	\\
				e^{\frac{t}{560}} \frac{dS}{dt} + \frac{1}{560} e^{\frac{t}{560}} \cdot S &=							\\
				\frac{d}{dt} \left[e^{\frac{t}{560}} \cdot S\right] &= e^{\frac{t}{560}} \cdot 0.05.
			\end{align*}
			Integrating this simplified equation with respect to $t$ provides
			\begin{align*}
				e^{\frac{t}{560}} \cdot S &= 0.05 \int e^{\frac{t}{560}} dt	\\
										  &= 28 e^{\frac{t}{560}} + C \quad \text{where $C$ is a constant},
			\end{align*}
			and finally reduces to 
			\begin{align*}
				S(t) = 28 + Ce^{\frac{-t}{560}}.
			\end{align*}
			It is given that the air is initially clean, that is at $t = 0$, $S(0) = 0$. Substituting this into $S(t)$ gives
			\begin{align*}
				S(0) &= 28 + Ce^{\frac{0}{560}} = 0	\\
				C &= -28.
			\end{align*}
			Therefore, $S(t)$ for $0 \leq t \leq 72$ is given by
			\begin{align*}
				S(t) &= 28 - 28e^{\frac{-t}{560}}	\\
					 &= 28 \left(1 - e^{\frac{-t}{560}}\right).
			\end{align*}
			\\
		\item What is the level of pollution in mg per cubic metre after 72 seconds.
			\\ \\
			The level of pollution in the room is given by $\frac{S(72)}{84}$ since the room contains 84 cubic metres of air. Hence, the level of pollution in the room after 72 seconds is
			\begin{align*}
				\frac{S(72)}{84} &= \frac{28 \cdot \left(1 - e^{\frac{-72}{560}}\right)}{84}	\\
								 &= \frac{1}{3} \cdot \left(1 - e^{\frac{-9}{70}}\right)		\\
								 &\approx 0.04021641686 \; \text{mg m$^{-3}$ (10 s.f.)}.
			\end{align*}
			\\
		\item How long does it take for the level of pollution to fall to 0.005 mg per cubic metre after the toast is taken outside.
			\\ \\
			After the toast is taken outside, smoke is no longer entering the room. Therefore, after 72 seconds the smoke leaving the room is governed by the equation
			\begin{align*}
				\frac{dS_n}{dt} = -\frac{S_n}{560}.
			\end{align*}
			This equation can be solved simply as
			\begin{align*}
				S_n(t) = S_0 e^{\frac{-t}{560}},
			\end{align*}
			where $S_0 = \frac{S(72)}{84} = \frac{1}{3} \cdot \left(1 - e^{\frac{-9}{70}}\right)$ and $t$ is time in seconds after the toast was removed. Solving for $\frac{S_n(t)}{84} = 0.005$ gives
			\begin{align*}
				\frac{S_n(t)}{84} &= 0.005	\\
				\frac{1}{3} \cdot \left(1 - e^{\frac{-9}{70}}\right) \cdot e^{\frac{-t}{560}} &= 0.005	\\
				e^{\frac{-t}{560}} &= \frac{3}{200 \cdot \left(1 - e^{\frac{-9}{70}}\right)}			\\
				\frac{-t}{560} &= \ln \left( \frac{3}{200 \cdot \left(1 - e^{\frac{-9}{70}}\right)} \right)	\\
				t &= -560 \cdot \ln\left( \frac{3}{200 \cdot \left(1 - e^{\frac{-9}{70}}\right)} \right)	\\
				  &\approx 1167.508933 \; \text{seconds (10 s.f.)}.
			\end{align*}
			Therefore it takes appropriately 1167.508933 seconds after the toast is removed from the room for the level of pollution to fall to 0.005 mg per cubic metre.
	\end{enumerate}
\end{document}
